% !TEX root = ../main.tex
%% Fiche 01 — Suites numériques (programme de Terminale spécialité, BO spécial n°8 du 25/07/2019)
%% Toutes les valeurs numériques sont vérifiées par verifications/fiche_01.py.

\fiche[suites]{1}{Suites numériques}
  {Récurrence\sepmot Limites\sepmot Comparaison~et~gendarmes\sepmot
   Convergence~monotone\sepmot Suites~$u_{n+1}=f(u_n)$\sepmot Algorithmes~de~seuil}

%% ---------------------------------------------------------------------
\sect{L'indispensable de Première}

{\otTableau
\begin{tabularx}{\linewidth}{@{}L{23mm}YY@{}}
 & \otTH{Suite arithmétique de raison $r$} & \otTH{Suite géométrique de raison $q$} \\
\midrule
\otCaps[8]{Relation} & $u_{n+1}=u_n+r$ & $u_{n+1}=q\,u_n$ \\
\otCaps[8]{Terme général} & $u_n=u_0+nr=u_p+(n-p)\,r$ & $u_n=u_0\,q^{n}=u_p\,q^{n-p}$ \\
\otCaps[8]{Somme} & $u_0+\dots+u_n=(n+1)\times\dfrac{u_0+u_n}{2}$
  & $u_0+\dots+u_n=u_0\times\dfrac{1-q^{n+1}}{1-q}\quad(q\neq1)$ \\
\otCaps[8]{Variations} & croissante si $r>0$, décroissante si $r<0$
  & si $u_0>0$ : croissante si $q>1$, décroissante si $0<q<1$ \\
\bottomrule
\end{tabularx}}

\vspace{1mm}
\noindent
\begin{minipage}[t]{0.485\linewidth}
\formule{1+2+\dots+n=\frac{n(n+1)}{2}}
\end{minipage}\hfill
\begin{minipage}[t]{0.485\linewidth}
\formule{1+q+\dots+q^{n}=\frac{1-q^{n+1}}{1-q}\quad(q\neq1)}
\end{minipage}

\begin{paire}
\begin{methode}[Étudier le sens de variation]
\begin{itemize}
  \item Étudier le signe de $u_{n+1}-u_n$.
  \item Termes tous strictement positifs : comparer $\frac{u_{n+1}}{u_n}$ à $1$.
  \item Si $u_n=f(n)$ : variations de $f$ sur $\intFO{0}{+\infty}$.
  \item Si $u_{n+1}=f(u_n)$ : par récurrence (§\,5).
\end{itemize}
\end{methode}
\begin{methode}[Suite $u_{n+1}=a\,u_n+b$, avec $a\neq1$]
\begin{etapes}
  \item Chercher $\ell$ tel que $\ell=a\ell+b$, soit $\ell=\frac{b}{1-a}$.
  \item Poser $v_n=u_n-\ell$ : $v_{n+1}=a\,v_n$, donc $(v_n)$ est géométrique.
  \item Conclure : $u_n=\ell+(u_0-\ell)\,a^{n}$.
\end{etapes}
\end{methode}
\end{paire}

%% ---------------------------------------------------------------------
\sect{Le raisonnement par récurrence}

\noindent
\begin{minipage}[c]{0.565\linewidth}
\begin{theoreme}[Principe de récurrence]
Soit $P(n)$ une propriété dépendant d'un entier naturel $n$, et $n_0\in\N$. Si
\begin{itemize}
  \item \etape{Initialisation} $P(n_0)$ est vraie,
  \item \etape{Hérédité} pour tout entier $k\geq n_0$, $P(k)$ vraie entraîne $P(k+1)$ vraie,
\end{itemize}
alors $P(n)$ est vraie pour tout entier $n\geq n_0$.
\end{theoreme}
\end{minipage}\hfill
\begin{minipage}[c]{0.405\linewidth}\centering
\begin{tikzpicture}[x=1mm,y=1mm,line join=round]
  \fill[otRule] (-2,-0.6) rectangle (64,0);
  \begin{scope}[shift={(4,0)},rotate=-28]
    \fill[otGreen,rounded corners=0.5pt] (0,0) rectangle (3,14);
  \end{scope}
  \foreach \i in {1,...,6}{\fill[otNavy,rounded corners=0.5pt] ({4+9*\i},0) rectangle ++(3,14);}
  \draw[otGreenDark,line width=0.6pt,-{Stealth[length=1.6mm]}] (1,18) to[bend left=35] (10.5,16.5);
  \node[font=\otSemi\fontsize{6}{7}\selectfont,text=otGreenDark,anchor=south] at (4,19.5) {INITIALISATION};
  \draw[otNavy,line width=0.6pt,-{Stealth[length=1.6mm]}] (32.5,16) to[bend left=40] (41.5,16);
  \node[font=\otSemi\fontsize{6}{7}\selectfont,text=otNavy,anchor=south] at (37,19.5) {HÉRÉDITÉ};
  \node[font=\scriptsize,text=otMuted,anchor=north] at (5.5,-1.2) {$n_0$};
  \node[font=\scriptsize,text=otMuted,anchor=north] at (32.5,-1.2) {$k$};
  \node[font=\scriptsize,text=otMuted,anchor=north] at (41.5,-1.2) {$k+1$};
\end{tikzpicture}\\[1mm]
\legende{Le premier domino tombe et chacun fait tomber le suivant : tous tombent.}
\end{minipage}

\begin{paire}
\begin{redaction}[le modèle à reproduire]
\fontsize{9.6}{12}\selectfont
Pour $n\in\N$, on note $P(n)$ : « $0\leq u_n\leq3$ ».\par
\etape{Initialisation} $u_0=0$, donc $P(0)$ est vraie.\par
\etape{Hérédité} Soit $k\in\N$ tel que $P(k)$ est vraie : $0\leq u_k\leq3$. Montrons $P(k+1)$.
[…calcul utilisant $0\leq u_k\leq3$…] Donc $0\leq u_{k+1}\leq3$.\par
\etape{Conclusion} Par récurrence, $P(n)$ est vraie pour tout $n\in\N$.
\end{redaction}
\begin{attention}
\begin{itemize}
  \item Sans initialisation, rien n'est prouvé : « $4^{n}+1$ est divisible par $3$ » est héréditaire,
    car \mbox{$4^{k+1}+1=4(4^{k}+1)-3$}, mais toujours fausse.
  \item Dans l'hérédité, on \emph{utilise} $P(k)$ pour un $k$ fixé ; on ne suppose jamais
    « $P(n)$ vraie pour tout $n$ ».
\end{itemize}
\end{attention}
\end{paire}

%% ---------------------------------------------------------------------
\sect{Suites bornées et limites}

\begin{definition}[Suites majorées, minorées, bornées]
$(u_n)$ est \cle{majorée} s'il existe un réel $M$ tel que $u_n\leq M$ pour tout $n$ ;
\cle{minorée} s'il existe $m$ tel que $u_n\geq m$ pour tout $n$ ; \cle{bornée} si elle est
majorée et minorée.
\end{definition}

\begin{paire}
\begin{definition}[Limite finie]
$(u_n)$ \cle{converge} vers $\ell$ si tout intervalle ouvert contenant $\ell$ contient tous les
termes $u_n$ à partir d'un certain rang. On note $\displaystyle\lim_{n\to+\infty}u_n=\ell$ ou
$u_n\to\ell$ ; cette limite est \cle{unique}.\par\vspace{1mm}
\centering
\begin{tikzpicture}
\begin{axis}[otSchema,width=80mm,height=33mm,xmin=0,xmax=15.5,ymin=0,ymax=3.6,
  xtick=\empty,ytick=\empty,xlabel={$n$},xlabel style={at={(1,0)},anchor=west}]
  \fill[otGreen,opacity=0.14] (axis cs:0,1.62) rectangle (axis cs:15.5,2.38);
  \draw[otGreen,line width=0.45pt] (axis cs:0,1.62) -- (axis cs:15.5,1.62)
    (axis cs:0,2.38) -- (axis cs:15.5,2.38);
  \draw[otNavy,dashed,line width=0.5pt] (axis cs:0,2) -- (axis cs:15.5,2);
  \draw[otMuted,densely dotted,line width=0.5pt] (axis cs:4.6,0) -- (axis cs:4.6,3.4);
  \addplot[only marks,mark=*,mark size=1.3pt,draw=otNavy,fill=white,samples at={0,...,4}]
    {2+1.4*cos(180*x)*0.75^x};
  \addplot[only marks,mark=*,mark size=1.3pt,otGreen,samples at={5,...,15}]
    {2+1.4*cos(180*x)*0.75^x};
  \node[anchor=east,font=\small,text=otNavy] at (axis cs:0,2) {$\ell$};
  \node[anchor=north,font=\otMedium\fontsize{6.5}{8}\selectfont,text=otMuted] at (axis cs:4.6,0) {rang $N$};
\end{axis}
\end{tikzpicture}\par
\legende{À partir du rang $N$, tous les termes restent dans la bande.}
\end{definition}
\begin{definition}[Limite infinie]
$u_n\to+\infty$ si tout intervalle $\intOO{A}{+\infty}$ contient tous les termes $u_n$ à partir
d'un certain rang. De même $u_n\to-\infty$ avec les intervalles $\intOO{-\infty}{A}$.\par\vspace{1mm}
\centering
\begin{tikzpicture}
\begin{axis}[otSchema,width=80mm,height=33mm,xmin=0,xmax=14.5,ymin=0,ymax=3.8,
  xtick=\empty,ytick=\empty,xlabel={$n$},xlabel style={at={(1,0)},anchor=west}]
  \fill[otGreen,opacity=0.14] (axis cs:0,1.8) rectangle (axis cs:14.5,3.75);
  \draw[otGreen,line width=0.45pt] (axis cs:0,1.8) -- (axis cs:14.5,1.8);
  \draw[otMuted,densely dotted,line width=0.5pt] (axis cs:6.6,0) -- (axis cs:6.6,3.6);
  \addplot[only marks,mark=*,mark size=1.3pt,draw=otNavy,fill=white,samples at={0,...,6}]
    {0.25+0.25*x+0.35*sin(deg(1.8*x))};
  \addplot[only marks,mark=*,mark size=1.3pt,otGreen,samples at={7,...,13}]
    {0.25+0.25*x+0.35*sin(deg(1.8*x))};
  \node[anchor=east,font=\small,text=otGreenDark] at (axis cs:0,1.8) {$A$};
  \node[anchor=north,font=\otMedium\fontsize{6.5}{8}\selectfont,text=otMuted] at (axis cs:6.6,0) {rang $N$};
\end{axis}
\end{tikzpicture}\par
\legende{Aussi grand que soit $A$, la suite finit par le dépasser définitivement.}
\end{definition}
\end{paire}

\noindent
\begin{minipage}[t]{0.385\linewidth}
\otIntertitre{Limites de référence}
{\otTableau
\begin{tabularx}{\linewidth}{@{}Yl@{}}
\toprule
$\frac1n$,\enspace $\frac1{n^{2}}$,\enspace $\frac1{\sqrt n}$,\enspace $\frac1{n^{p}}$ & $\to0$ \\
$n$,\enspace $n^{2}$,\enspace $\sqrt n$,\enspace $n^{p}$\enspace($p\geq1$) & $\to+\infty$ \\
\midrule
$q>1$ & $q^{n}\to+\infty$ \\
$q=1$ & $q^{n}=1$ \\
$-1<q<1$ & $q^{n}\to0$ \\
$q\leq-1$ & pas de limite \\
\bottomrule
\end{tabularx}}
\end{minipage}\hfill
\begin{minipage}[t]{0.575\linewidth}
\otIntertitre{Opérations sur les limites ($\ell$, $\ell'$ réels)}
{\otTableau\fontsize{9.4}{11}\selectfont
\begin{tabularx}{\linewidth}{@{}L{13mm}CCCC@{}}
\toprule
 & $\ell$ et $\ell'$ & $\ell$ et $+\infty$ & $+\infty$ et $+\infty$ & $+\infty$ et $-\infty$ \\
\midrule
\otCaps[6]{Somme} & $\ell+\ell'$ & $+\infty$ & $+\infty$ & \cellcolor{otOrangeTint}\textbf{F.I.} \\
\otCaps[6]{Produit} & $\ell\ell'$ & $\pm\infty$\,\textsuperscript{(1)} & $+\infty$ & $-\infty$ \\
\bottomrule
\end{tabularx}}\par\vspace{1mm}
{\fontsize{9}{11.4}\selectfont\RaggedRight
\textsuperscript{(1)} Règle des signes si $\ell\neq0$ ; \colorbox{otOrangeTint}{\textbf{F.I.}} si $\ell=0$.\\
Quotient : $\frac{\ell}{\pm\infty}\to0$ ; $\frac{\ell}{0}$ avec $\ell\neq0$ donne $\pm\infty$
(signe du dénominateur).\par}
\vspace{1.2mm}
{\RaggedRight Les quatre \cle{formes indéterminées} :\quad
$\infty-\infty$,\quad $0\times\infty$,\quad $\dfrac{\infty}{\infty}$,\quad $\dfrac00$.\par}
\end{minipage}

\begin{methode}[Lever une forme indéterminée]
{\otTableau\renewcommand{\arraystretch}{1.6}%
\begin{tabularx}{\linewidth}{@{}L{25mm}>{\RaggedRight\arraybackslash}p{47mm}Y@{}}
\otCaps[6]{Quotient de polynômes} & factoriser par le terme de plus haut degré
  & $\dfrac{3n^{2}-n}{1-2n^{2}}=\dfrac{3-\frac1n}{\frac1{n^{2}}-2}\to\res{-\dfrac32}$ \\
\otCaps[6]{Différence de racines} & multiplier par la quantité conjuguée
  & $\sqrt{n+1}-\sqrt n=\dfrac{1}{\sqrt{n+1}+\sqrt n}\to\res{0}$ \\
\otCaps[6]{Puissances} & factoriser par la puissance dominante
  & $3^{n}-2^{n}=3^{n}\Bigl(1-\bigl(\tfrac23\bigr)^{n}\Bigr)\to\res{+\infty}$ \\
\end{tabularx}}
\end{methode}

\begin{attention}[les pièges classiques]
\begin{itemize}
  \item $\infty-\infty$ et $0\times\infty$ ne valent ni $0$ ni $1$ : ce sont des formes indéterminées,
    il faut transformer l'expression.
  \item $\bigl((-1)^{n}\bigr)$ est bornée mais n'a pas de limite : une suite bornée ne converge pas forcément.
  \item Pour $(q^{n})$, tout dépend de $q$ : $0{,}9^{n}\to0$ mais $1{,}1^{n}\to+\infty$.
  \item On n'écrit « $\lim u_n=\dots$ » qu'après avoir justifié que cette limite existe.
\end{itemize}
\end{attention}

%% ---------------------------------------------------------------------
\sect{Comparaison et convergence}

\noindent
\begin{minipage}[c]{0.585\linewidth}
\begin{theoreme}[Comparaison]
Si $u_n\leq v_n$ à partir d'un certain rang et $u_n\to+\infty$, alors $v_n\to+\infty$.\par
Si $v_n\leq u_n$ à partir d'un certain rang et $u_n\to-\infty$, alors $v_n\to-\infty$.
\end{theoreme}
\begin{theoreme}[Théorème des gendarmes]
Si $u_n\leq v_n\leq w_n$ à partir d'un certain rang, et si $(u_n)$ et $(w_n)$ convergent vers
\emph{la même} limite $\ell$, alors $(v_n)$ converge vers $\ell$.\par\vspace{0.6mm}
{\fontsize{9.4}{11.6}\selectfont Exemple : $-\frac1n\leq\frac{\sin n}{n}\leq\frac1n$, donc
$\frac{\sin n}{n}\to0$.}
\end{theoreme}
\end{minipage}\hfill
\begin{minipage}[c]{0.385\linewidth}\centering
\begin{tikzpicture}
\begin{axis}[otSchema,width=68mm,height=46mm,xmin=0,xmax=15,ymin=0.4,ymax=3.6,
  xtick=\empty,ytick=\empty,xlabel={$n$},xlabel style={at={(1,0)},anchor=west}]
  \draw[otNavy,dashed,line width=0.5pt] (axis cs:0,2) -- (axis cs:15,2);
  \addplot[only marks,mark=*,mark size=1.2pt,draw=otNavy,fill=white,samples at={1,...,14}] {2+1.5/(x+0.5)};
  \addplot[only marks,mark=*,mark size=1.2pt,draw=otNavy,fill=white,samples at={1,...,14}] {2-1.5/(x+0.5)};
  \addplot[only marks,mark=*,mark size=1.5pt,otGreen,samples at={1,...,14}] {2+1.5*sin(deg(1.3*x))/(x+0.5)};
  \node[anchor=east,font=\small,text=otNavy] at (axis cs:0,2) {$\ell$};
  \node[anchor=west,font=\scriptsize,text=otNavy] at (axis cs:14.3,2.28) {$w_n$};
  \node[anchor=west,font=\scriptsize,text=otNavy] at (axis cs:14.3,1.72) {$u_n$};
  \node[anchor=south,font=\scriptsize,text=otGreenDark] at (axis cs:3,3.02) {$v_n$};
\end{axis}
\end{tikzpicture}\\
\legende{Prise en étau entre $(u_n)$ et $(w_n)$, la suite $(v_n)$ converge vers $\ell$.}
\end{minipage}

\begin{paire}
\begin{theoreme}[Convergence monotone — admis]
Toute suite \cle{croissante et majorée} converge.\par
Toute suite \cle{décroissante et minorée} converge.\par\vspace{1mm}
\otEtiquette{otNavy}{Conséquences}\par\vspace{0.6mm}
Une suite croissante \emph{non majorée} tend vers $+\infty$.\par
Si $(u_n)$ croît et converge vers $\ell$, alors $u_n\leq\ell$ pour tout $n$.
\end{theoreme}
\begin{attention}
\begin{itemize}
  \item Si la suite est croissante et majorée par $M$, sa limite vérifie \mbox{$\ell\leq M$}
    sans être forcément égale à $M$ : $1-\frac{1}{n+1}$ est majorée par $5$ et tend vers $1$.
  \item Le théorème donne l'\emph{existence} de la limite, pas sa valeur.
  \item Gendarmes : les deux suites qui encadrent doivent avoir la \emph{même} limite.
\end{itemize}
\end{attention}
\end{paire}

\begin{propriete}[Lien avec les fonctions]
Si $u_n=f(n)$ et $\displaystyle\lim_{x\to+\infty}f(x)=L$, alors $u_n\to L$. Si $u_n\to\ell$ et si
$f$ est continue en $\ell$, alors $f(u_n)\to f(\ell)$.
\end{propriete}

\begin{demonstration}[les trois preuves exigibles]
\etape{1. Croissante non majorée}
Soit $A$ réel : $A$ ne majore pas la suite, il existe $N$ tel que $u_N>A$ ; comme $(u_n)$ croît,
$u_n\geq u_N>A$ pour tout $n\geq N$. Donc $u_n\to+\infty$.\par\vspace{0.6mm}
\etape{2. Minorée par une suite de limite $+\infty$}
Si $v_n\geq u_n$ pour $n\geq n_1$ et $u_n\to+\infty$ : pour tout $A$, il existe $n_2$ tel que
$u_n>A$ pour $n\geq n_2$ ; alors $v_n>A$ pour $n\geq\max(n_1,n_2)$.\par\vspace{0.6mm}
\etape{3. Limite de $(q^{n})$}
Par récurrence, $(1+a)^{n}\geq1+na$ pour $a>0$ (\emph{inégalité de Bernoulli}). Si $q>1$, avec
$a=q-1$ : $q^{n}\geq1+n(q-1)\to+\infty$, donc $q^{n}\to+\infty$ par comparaison. Si $0<|q|<1$,
on applique ce résultat à $\frac{1}{|q|}>1$ : $|q|^{n}\to0$.
\end{demonstration}

%% ---------------------------------------------------------------------
\sect{Suites définies par une relation $u_{n+1}=f(u_n)$}

\noindent
\begin{minipage}[c]{0.585\linewidth}
\begin{methode}[Étude complète, en quatre temps]
\begin{etapes}
  \item \cle{Encadrer} : montrer par récurrence que $u_n\in I$ pour tout $n$, avec
    $f(I)\subset I$ (souvent grâce aux variations de $f$).
  \item \cle{Sens de variation} : par récurrence (si $f$ est croissante sur $I$, comparer $u_0$
    et $u_1$), ou par le signe de $u_{n+1}-u_n$.
  \item \cle{Convergence} : la suite est monotone et bornée, donc elle converge.
  \item \cle{Limite} : si $u_n\to\ell$ et $f$ est continue en $\ell$, alors $\ell=f(\ell)$ ; on
    résout cette équation et on garde la solution compatible avec l'encadrement.
\end{etapes}
\end{methode}
\begin{exemple}[$u_0=0$ et $u_{n+1}=\sqrt{u_n+6}$]
Par récurrence, $0\leq u_n\leq u_{n+1}\leq3$ : la suite converge vers $\ell\geq0$ tel que
$\ell=\sqrt{\ell+6}$, soit $\ell^{2}-\ell-6=0$. Donc $\ell=\res{3}$ (la racine $-2$ est exclue).
\end{exemple}
\end{minipage}\hfill
\begin{minipage}[c]{0.385\linewidth}\centering
\begin{tikzpicture}
\begin{axis}[otSchema,width=66mm,height=62mm,xmin=0,xmax=4,ymin=0,ymax=4,
  xtick=\empty,ytick={3},xlabel={$x$},ylabel={$y$},
  xlabel style={at={(1,0)},anchor=west},ylabel style={at={(0,1)},anchor=south,rotate=-90}]
  \addplot[otRule!80!otMuted,line width=0.6pt,domain=0:3.9] {x};
  \addplot[otNavy,line width=1.1pt,domain=0:3.9,samples=60] {sqrt(x+6)};
  \addplot[otGreen,line width=0.8pt] coordinates
    {(0,0) (0,2.4495) (2.4495,2.4495) (2.4495,2.9068) (2.9068,2.9068) (2.9068,2.9844) (2.9844,2.9844) (2.9844,2.9974)};
  \draw[otGreen!70!otMuted,densely dotted,line width=0.5pt] (axis cs:2.4495,2.4495) -- (axis cs:2.4495,0);
  \draw[otGreen!70!otMuted,densely dotted,line width=0.5pt] (axis cs:2.9068,2.9068) -- (axis cs:2.9068,0);
  \draw[otNavy!60,densely dotted,line width=0.5pt] (axis cs:0,3) -- (axis cs:3,3);
  \addplot[only marks,mark=*,mark size=1.2pt,otGreen] coordinates {(0,0) (2.4495,0) (2.9068,0)};
  \fill[otNavy] (axis cs:3,3) circle (1.5pt);
  \node[anchor=north,font=\scriptsize] at (axis cs:0.08,-0.05) {$u_0$};
  \node[anchor=north east,font=\scriptsize] at (axis cs:2.5,-0.05) {$u_1$};
  \node[anchor=north west,font=\scriptsize] at (axis cs:2.86,-0.05) {$u_2$};
  \node[anchor=south,font=\scriptsize,text=otNavy] at (axis cs:1,2.9) {$y=\sqrt{x+6}$};
  \node[anchor=west,font=\scriptsize,text=otMuted] at (axis cs:3.55,3.55) {$y=x$};
\end{axis}
\end{tikzpicture}\\
\legende{La « toile d'araignée » : on monte à la courbe, on revient à la droite $y=x$ ; les
termes s'accumulent vers le point fixe $3$.}
\end{minipage}

%% ---------------------------------------------------------------------
\sect{Algorithmes à connaître}

\noindent
\begin{minipage}[t]{0.485\linewidth}
\otIntertitre{Recherche d'un seuil}
\begin{python}
def seuil(s):
    n, u = 0, 0         # u0 = 0
    while u <= s:
        u = 0.5 * u + 1
        n = n + 1
    return n
\end{python}
{\fontsize{9.4}{11.8}\selectfont\RaggedRight Pour $u_0=0$ et $u_{n+1}=0{,}5\,u_n+1$, on a
$u_n=2-2\times0{,}5^{n}$. \texttt{seuil(1.99)} renvoie $\res{8}$ : premier rang où $u_n>1{,}99$
($u_7\approx1{,}984$ et $u_8\approx1{,}992$). La boucle s'arrête car $u_n\to2>1{,}99$.\par}
\end{minipage}\hfill
\begin{minipage}[t]{0.485\linewidth}
\otIntertitre{Somme des premiers termes}
\begin{python}
def somme(n):
    u, S = 1, 1         # u0 = 1
    for k in range(n):
        u = 3 * u
        S = S + u
    return S            # u0 + ... + un
\end{python}
{\fontsize{9.4}{11.8}\selectfont\RaggedRight \texttt{somme(4)} renvoie
$1+3+9+27+81=\res{121}=\frac{3^{5}-1}{3-1}$. De même, $u_0=1$ et
$u_{n+1}=\frac12\bigl(u_n+\frac{2}{u_n}\bigr)$ donnent des valeurs approchées de $\sqrt2$ :
$1{,}5$ ; $1{,}416\,7$ ; $1{,}414\,22$\,…\par}
\end{minipage}

%% ---------------------------------------------------------------------
\strategie{Comment trouver la limite d'une suite ?}{%
  {$u_n$ est donnée en fonction de $n$}/{Limites de référence et opérations ; en cas de forme
    indéterminée : factoriser, quantité conjuguée.},
  {On dispose d'un encadrement ou d'une inégalité}/{Gendarmes (limite finie) ou théorème de
    comparaison (limite infinie).},
  {$u_{n+1}=a\,u_n+b$ avec $a\neq1$}/{Suite auxiliaire $v_n=u_n-\ell$, géométrique, puis limite
    de $q^{n}$.},
  {$u_{n+1}=f(u_n)$}/{Monotone et bornée, donc convergente ; puis $\ell=f(\ell)$ si $f$ est
    continue en $\ell$.},
  {On demande seulement si elle converge}/{Convergence monotone : croissante et majorée, ou
    décroissante et minorée.}}

\begin{jesais}
  \item calculer un terme et une somme (arithmétique, géométrique) ;
  \item étudier le sens de variation d'une suite ;
  \item rédiger une récurrence complète, initialisation comprise ;
  \item calculer une limite et lever une forme indéterminée ;
  \item utiliser la comparaison et les gendarmes ;
  \item justifier une convergence par la convergence monotone ;
  \item étudier une suite $u_{n+1}=f(u_n)$ et trouver sa limite ;
  \item écrire un algorithme de seuil en Python.
\end{jesais}

\pourAllerPlusLoin{Livre OptiTAB \emph{Mathématiques Terminale}, chapitre 1 « Suites
numériques » : 15 exercices corrigés pas à pas, des applications directes aux problèmes type bac.
Cours, quiz et exercices en ligne sur \textbf{\otSite}.}
